\section{B. The mixed topos and the cohomology of hole filling}\label{part:B}

We first fix a consequential ambiguity. The assertions about cohomology on
LM scans 127--129 concern an ordinary glued topos. Deriving the gluing
functor is indispensable. Its underived value does not remember the kernel
of a local monodromy homomorphism. This is precisely where the proposed
proof on scan 129 can fail.

\subsection{The actual construction}

\paragraph{B.1 (Proved): the Artin gluing of LM \S14.}\label{B1}
\textit{Manuscript: LM \S14, folio 60--62, scans 123--127.}
Let $G$, $H_s$, and $Q_s$ be discrete groups, or profinite groups with
continuous actions on discrete sets. Let
\[
 r_s:H_s\longrightarrow G,\qquad q_s:H_s\twoheadrightarrow Q_s,
 \qquad I_s=\ker q_s.
\]
We use a finite index set $S$ for the cohomological assertions below;
the construction itself permits an arbitrary small index set. There is a
Grothendieck topos $\mathcal B$ whose objects and morphisms are
\[
 (M,(N_s), (a_s)),\qquad
 a_s:q_s^*N_s\longrightarrow r_s^*M,
\]
and componentwise equivariant maps commuting with the $a_s$.
Equivalently, $a_s$ is a $Q_s$-equivariant map
$N_s\to(r_s^*M)^{I_s}$. The functor of gluing is
$L=(q_{s*}r_s^*)_s$. Its open part is $\beta G$, its complementary
closed part is $\coprod_s\beta Q_s$, and it is the topos denoted
$B_{X,U}$, $B_{U,X}$, and subsequently $B_{X;S}$ in LM
\S14, folios 60--62, scans 123--127, and \S16, scans 163--166.

\begin{proof}
The functor $L$ is accessible and preserves finite limits: restriction
does so, and fixed points are a right adjoint. The Artin gluing theorem
\citep[IV, Th\'eor\`eme 9.5.4]{SGA4} applies. Here is also the concrete
identification needed for the manuscript. In the displayed adjoint
description, finite limits and arbitrary colimits are componentwise,
since $r_s^*$ and $q_s^*$ preserve both. The subterminal
$u=(1,(\varnothing),a)$ has slice $\mathcal B/u\simeq\beta G$:
an object mapping to $u$ has all $N_s$ empty. The complementary closed
embedding sends $(N_s)$ to $(1,(N_s),!)$. The right adjoint to restriction
to the open part is $M\mapsto(M,L(M),\mathrm{id})$, so restriction of
this right adjoint to the closed part is exactly $L$, as asserted on
scan 127. This also proves that the closed embedding is not generally an
open induction morphism. Evaluation on the open and closed fibres is a
conservative family of points. Every point belongs to one of these two
parts, because its inverse image sends $u$ to either $0$ or $1$.
\end{proof}

For a smooth connected curve $X/k$, a finite reduced closed subset
$S$, and $U=X-S$, take $G=\pi_1(U)$,
$H_s=\operatorname{Gal}(K_s^{\mathrm{sep}}/K_s)$ for the fraction field
of the henselization at $s$, and $Q_s=\operatorname{Gal}(k(s)^{\mathrm{sep}}/k(s))$.
The maps require compatible geometric points; replacing these choices
conjugates the maps and gives canonically equivalent gluing categories.
Artin recollement on $X_{\et}$ gives a fully faithful inverse-image
functor
\[
 f^*:\mathcal B\longrightarrow X_{\et},\qquad
 f:X_{\et}\longrightarrow\mathcal B.
\]
Its essential image consists of the sheaves whose restriction to $U$
is a continuous $G$-set. Infinite such objects are allowed; ``continuous''
is not being replaced by ``one finite quotient acts on the entire set''.
The henselian local topos itself is the gluing of $\beta H_s$ and
$\beta Q_s$ by $q_{s*}$: both its generic and closed strata are spectra
of fields. Thus the local equivalence asserted at the top of scan 129
is correct. The subsequent derived base-change inference is a different
assertion.

\paragraph{B.2 (Proved): filling holes at level one.}\label{B2}
\textit{Manuscript: LM \S14, folio 60, scans 123.}
In the preceding geometric setting the continuous surjection
\[
 \pi_1(U)\longrightarrow\pi_1(X)
\]
has kernel the \emph{closed} normal subgroup generated by the images
$r_s(I_s)$. For a discrete surface the same formula holds with ordinary
normal closure. This proves the formula of LM \S14, folio 60, scan 123.

\begin{proof}
A finite $G$-set extends to a finite locally constant sheaf on $X$
exactly when every $a_s$ can be an isomorphism. Equivalently $I_s$
acts trivially on its restriction. Indeed the finite cover of $U$ has
a unique normalization over the regular curve $X$, and at each missing
discrete valuation its normalization is \`etale precisely when inertia
acts trivially. The associated $Q_s$-set is then forced. Thus restriction
identifies finite covers of $X$ with the full subcategory of finite
$G$-sets annihilated by the asserted normal subgroup. This is exactly
the finite representation category of the quotient; the fibre functor
identifies its profinite automorphism group with $\pi_1(X)$. The proof
does not select a connected component of a cover and hence applies
componentwise to groupoids without adding base points to the conclusion.
\end{proof}

\paragraph{B.3 (Proved): the effective local group and the derived comparison.}\label{B3}
\textit{Manuscript: LM \S14, folio 62--63, scans 127--129.}
Write
\[
 K_s=\operatorname{im}(H_s\xrightarrow{(r_s,q_s)}G\times Q_s).
\]
In the profinite case this image is closed. For a torsion abelian object
$F=(M,N_s,a_s)$ of $\mathcal B$ there is a natural fibre sequence
\begin{equation}\label{eq:B-effective}
 R\Gamma(\mathcal B,F)\longrightarrow
 R\Gamma(G,M)\oplus\prod_sR\Gamma(Q_s,N_s)
 \longrightarrow\prod_sR\Gamma(K_s,M),
\end{equation}
where $K_s$ acts on $M$ through its projection to $G$, and the second
map is the difference of restriction and the map induced by $a_s$.
For the geometric curve, the corresponding formula is
\begin{equation}\label{eq:B-geometric}
 R\Gamma(X,f^*F)\longrightarrow
 R\Gamma(U,M)\oplus\prod_sR\Gamma(Q_s,N_s)
 \longrightarrow\prod_sR\Gamma(H_s,M).
\end{equation}
Consequently, whenever $U$ has the $K(\pi,1)$ comparison for $M$,
\begin{equation}\label{eq:B-defect}
 \operatorname{Cone}\bigl(R\Gamma(\mathcal B,F)\to R\Gamma(X,f^*F)\bigr)
 \simeq
 \prod_s\operatorname{Cone}\bigl(R\Gamma(K_s,M)\to R\Gamma(H_s,M)\bigr)[-1].
\end{equation}
This is an exact coefficientwise necessary and sufficient criterion,
not merely a sufficient injectivity hypothesis. In particular an
injective $(r_s,q_s)$ for each $s$ suffices.

\begin{proof}
For discrete groups form the category with one object $u$ with group
$G$, one object $s$ with group $Q_s$, no arrows from $u$ to $s$ or
between distinct closed objects, and crossing biset
$(G\times Q_s)/H_s$. We use the inverse on the second factor to make
this a left $G$, right $Q_s$ biset. Its representations are exactly
the triples of B.1. The action groupoid of its crossing arrows has
stabilizer $K_s$, not $H_s$: the kernel of $(r_s,q_s)$ acts trivially
on every crossing arrow. Thus this groupoid is equivalent to $BK_s$.

For clarity, the cochain assertion follows directly from the bar
resolution of this category. A nondegenerate chain lies entirely in
one stratum, or contains exactly one crossing arrow. The subcomplexes
of the first kind give the two stratum cochain complexes. The chains
of the second kind have automorphism strings on the two sides of the
crossing arrow; their two-sided bar complex is the cochain complex of
the crossing action groupoid. Its two augmentations are restriction
from $G$ and the specialization map from $Q_s$. The extra crossing
arrow shifts the latter complex by one. Taking the total complex is
therefore the homotopy fibre displayed in \eqref{eq:B-effective}.
This argument includes nonconstant coefficients: their value at a
crossing arrow is the value $M$ at its target, and its source value
maps there by $a_s$.

For profinite groups the finite continuous triples form a coherent
generating pretopos: they are closed under finite limits, finite
coproducts and quotients, and every continuous triple is a filtered
union of finite subtriples. Every such finite triple and every finite
diagram of them descends to a common finite quotient of $G$ and the
$Q_s$ and the corresponding finite images of the $H_s$. Thus our
topos is the coherent inverse limit of these finite-category topoi.
Their transition inverse images preserve finite objects and are
coherent. The cohomology continuity theorem
\citep[VI, Corollary~8.7.7, with 8.7.1 and Remark~8.7.8]{SGA4II}
therefore identifies its derived cohomology with the filtered
colimit of the finite-level computations. This verifies that the
following continuous bar calculation really computes topos
cohomology, rather than introducing a different cohomology theory.
Each continuous map
from the compact space of a fixed finite string of arrows to a
discrete torsion coefficient module has finite image and factors
through a finite quotient. Finite collections of such maps admit a
common quotient. Hence the continuous bar complex is the filtered
colimit of the finite bar complexes, with coefficients first restricted
to their finite submodules. Filtered colimits of abelian groups are
exact. The discrete argument therefore proves the same formula for
continuous cohomology.

For \eqref{eq:B-geometric}, use the localization triangle for $U\subset X$
and excision at the finitely many points of $S$. The stalk of
$Rj_*M$ at a geometric $s$ is $R\Gamma(I_s,M)$, with its residue
Galois action. The henselian local topos has derived sections
$R\Gamma(Q_s,N_s)$; after strict henselization this is the exact
closed-stalk functor, and descent gives the displayed expression.
Hochschild--Serre identifies its local punctured term with
$R\Gamma(H_s,M)$. This gives \eqref{eq:B-geometric}. The comparison
of the two fibre sequences is identity on the middle term when $U$
has the stipulated comparison. The elementary cone identity in a
stable category gives \eqref{eq:B-defect}.
\end{proof}

\paragraph{B.4 (Refuted): the proposition of scan 129 without a local hypothesis.}\label{B4}
\textit{Manuscript: LM \S14, folio 63, scans 129.}
Let $k$ be algebraically closed of characteristic zero,
$X=\mathbf P^1_k$, $S=\{\infty\}$, and $U=\mathbf A^1_k$.
Then $G=Q_s=1$, $H_s=I_s=\widehat{\mathbb Z}(1)$, and
$\mathcal B\simeq\operatorname{Fun}([1],\Set)$. For
$A=\mathbb Z/n$, $n>1$,
\[
 H^2(\mathcal B,A)=0,\qquad
 H^2(X_{\et},A)=A(-1)\ne0.
\]
Nevertheless $U$ is a $K(\pi,1)$. Thus the proposition on LM \S14,
folio 63, scan 129, and the prompt's inferred ``all nonempty $S$''
conclusion are false.

\begin{proof}
The triple is just a map $N\to M$. Its global sections are $N$,
because a compatible pair is uniquely specified by its source value.
This functor is exact on abelian objects, so all higher cohomology
vanishes. The target computation is the Kummer calculation of
\citep[Tag 03RQ, Lemma 59.69.1]{StacksCurveKummer}; it can also be read
as the cokernel of multiplication by $n$ on
$\operatorname{Pic}(\mathbf P^1)=\mathbb Z$.
In B.3, $K_s=1$, whereas $H_s=\widehat{\mathbb Z}(1)$.
The missing cone is $A(-1)[-2]$. In particular the error on scan 128
is the assertion that local-support derived functors can be computed
after restricting an injective global mixed object to the local mixed
topos. That restriction need not preserve injectives. Equivalence of
the local topoi does not prove this base-change assertion.
The program \texttt{code/root\_checks.py} independently computes the
Betti numbers of the nerve $[1]$ and a triangulated $S^2$ modulo
$2$, $3$, and $5$; the degree-two discrepancy is recorded in its output.
\end{proof}

\subsection{The complete curve comparison}

\paragraph{B.5 (Known, with reduction proved): the curve $K(\pi,1)$ theorem.}\label{B5}
\textit{Manuscript: LM \S14, folio 63, scans 129.}
For a smooth connected curve over any field, an affine curve and a
proper curve of geometric genus at least one are $K(\pi,1)$ for all
torsion locally constant abelian coefficients. The same holds for
their nonempty open subschemes. A proper curve of geometric genus
zero has the comparison for $p$-primary coefficients in
characteristic $p$, but not for every prime-to-$p$ coefficient.

\begin{proof}
The precise all-characteristic theorem is
\citet[Lemma 2.7(a)]{SchmidtStix2016}; their definition and the
cohomological criterion immediately preceding that lemma are the
ones needed here. We spell out the reduction to the present version.
Over an algebraic closure, an affine smooth curve has torsion
cohomological dimension at most one. For a finite local system,
first trivialize its monodromy on a finite \`etale cover. Every
degree-one class is a torsor under its finite abelian coefficient
group and is killed by that finite cover; higher classes vanish.
The effaceability criterion for the universal covering pro-object
therefore identifies its cohomology with continuous group cohomology.

For a proper curve of positive genus, after the same trivialization
only degree two needs attention. For $\ell\ne p$ the trace
identification is $H^2(Y,\mathbb Z/\ell^a(1))=\mathbb Z/\ell^a$;
pullback under a finite \`etale map is multiplication by its degree.
Such a curve has cyclic \`etale covers of degree $\ell^a$, by its
Jacobian or the prime-to-$p$ surface-group presentation. Its degree-two
class is therefore effaced. For $p$-primary coefficients in positive
characteristic, degree two already vanishes; this follows from
Artin--Schreier and the surjectivity of $F-1$ on $H^1(Y,\mathcal O_Y)$,
and extension through finite filtrations. These are the vanishing
statements in \citep[Tags 0A3J and 03SB]{StacksAS,StacksCurves}.
The same argument proves the $p$-primary assertion for genus zero.

For an arbitrary ground field, first replace it, if necessary, by
its finite separable algebraic closure in the curve's function field.
The curve is geometrically connected over this constant field;
neither its etale topos nor its fundamental group changes in this
rephrasing. Apply the compatible
Hochschild--Serre spectral sequences for the geometric curve and its
fundamental group extension by $\operatorname{Gal}(k^{\mathrm{sep}}/k)$.
The geometric comparisons just proved identify their second pages
and all differentials, hence their abutments. Purely inseparable
extensions do not change the \`etale site. Infinite discrete torsion
modules are filtered unions of finite invariant submodules: the orbit
of each element is finite by continuity, and the subgroup generated
by finitely many torsion elements of bounded order is finite.
Both cohomology theories commute with these filtered colimits for
the quasi-compact curves in question. This proves the assertion in
exactly the coefficient generality used in B.3.
\end{proof}

\paragraph{B.6 (Proved): globalization of local covering germs on curves.}\label{B6}
\textit{Manuscript: LM \S14, folio 63, scans 129.}
Let $U$ be a smooth connected affine curve over an algebraically
closed field $k$, let $\overline U$ be its smooth proper completion,
and let $s\in\overline U-U$. The full local inertia homomorphism
\[
 \operatorname{Gal}(k((t))^{\mathrm{sep}}/k((t)))
 \longrightarrow\pi_1(U)
\]
is injective except when $\operatorname{char}k=0$ and
$U\simeq\mathbf A^1$. In that exception it is the zero map from
$\widehat{\mathbb Z}(1)$ to $1$. The analogous tame local
homomorphism is injective, in every characteristic, precisely when
$U\not\simeq\mathbf A^1$.

\begin{proof}
We prove the full assertion by dominating every finite local covering
germ by the germ of a global finite \`etale covering of $U$.
For a continuous map $h:H\to G$ of profinite groups this criterion
implies injectivity: for each open normal $N\lhd H$, dominance gives
an open subgroup of $G$ whose inverse image in $H$ is contained in
$N$; the intersection of all such $N$ is $1$. Conversely, a closed
embedding $H\subset G$ has the criterion. Indeed, for an open
subgroup $V\subset H$ choose open normal $N\subset G$ with
$N\cap H\subset V$. Then $VN$ is an open subgroup of $G$ and
$VN\cap H=V$, so the $H$-orbit of its distinguished coset in
$G/VN$ is $H/V$. This also specifies the meaning of ``sub-germ''.

We next globalize cyclic local steps. For $n$ prime to the
characteristic, the local cyclic extension of degree $n$ is unique:
Hensel's lemma and divisibility of $k^*$ make every local unit an
$n$th power, so its Kummer class is determined by valuation modulo
$n$. If there are two boundary points $s,t$, divisibility of
$\operatorname{Pic}^0(\overline U)$ supplies a divisor $E$ with
$n[E]=[t-s]$. Thus $s-t+nE$ is principal, and its Kummer covering is
\`etale on $U$ and has ramification index $n$ at $s$.

If there is one boundary point and the genus is positive, use
\citep[XIII, Corollaire 2.12]{SGA1}. In its prime-to-$p$ presentation,
the peripheral generator is the inverse of the product of
commutators of $2g$ handle generators. Map one handle pair to
\[
 x=(1,0,0),\quad y=(0,1,0)
 \quad\text{in }\operatorname{UT}_3(\mathbb Z/n),
\]
where multiplication on triples is
$(a,b,c)(d,e,f)=(a+d,b+e,c+f+ae)$. Map all other handle generators
to $1$. The commutator is $(0,0,1)$ of order $n$. The cited
existence clause gives a tame cover with local ramification index
$n$ at the unique cusp. Its global group need not be cyclic, which
is essential here. The orders for $n=2,3,4,5,9$ are checked by the
attached program; the displayed multiplication proves the result
for every $n$.

In characteristic $p$, every cyclic local extension of degree $p$
is Artin--Schreier. It can be globalized without adding a branch
point inside $U$. Put $P=\bigoplus_{x\in\overline U-U}
k((t_x))/k[[t_x]]$. Principal parts give an exact sequence
\[
 \Gamma(U,\mathcal O_U)\longrightarrow P
 \xrightarrow{\delta} H^1(\overline U,\mathcal O_{\overline U})
 \longrightarrow0.
\]
The map is Frobenius compatible. The map $F-1$ on the last,
finite-dimensional, $k$-vector space is surjective
\citep[Tag 0A3J, Lemma 59.63.2]{StacksAS}. Given a desired local
principal part $a$ at $s$, with zero parts at the other boundary
points, choose $b\in P$ with
$(F-1)\delta(b)=\delta(a)$. Then $a-(F-1)b$ is the principal part
of a function $f$ regular on $U$. The covering $z^p-z=f$ has the
required local class. The integral remainder is killed by
Artin--Schreier on the henselian local ring, since its residue
field is algebraically closed. Thus the principal-part calculation
really determines the finite local covering germ.

There is just one obstruction to starting this procedure:
the one-cusp genus-zero curve has no nontrivial tame cover. In
characteristic $p$ first replace it by the finite \`etale cover
$y^p-y=x^m$, where $p\nmid m$ and
$(p-1)(m-1)/2>0$. Its completion has that positive genus and one
point above infinity. The genus follows from Riemann--Hurwitz:
the pole of order $m$ has different exponent $(p-1)(m+1)$, hence
$2g-2=-2p+(p-1)(m+1)$. The different formula follows directly
from the single ramification break $m$ of this Artin--Schreier
extension. We use $(p,m)=(2,3),(3,2)$ for the small characteristics.

Every finite local Galois group over $k((t))$ is an extension of
a $p$-group by a cyclic group of order prime to $p$, hence is
solvable. Choose a tower with cyclic prime-degree steps. Starting
with the preceding positive-genus cover when necessary, suppose
a global completion contains the $j$th local field in the tower.
The compositum with the next field is cyclic of that prime degree
or trivial. The cyclic construction just given dominates it by a
new global cover. Induction on the fixed finite tower terminates.
If the initial auxiliary cover was not contained in the prescribed
extension, take composita throughout; cyclicity survives this base
change. Subsequent finite separable covers preserve positive genus
by Riemann--Hurwitz, or preserve the existence of at least two
boundary points. The construction therefore never loses its
applicability. This proves injectivity in positive characteristic,
and also in characteristic zero outside $\mathbf A^1$.

For tame inertia it suffices to detect each finite cyclic quotient
directly using the Kummer or unitriangular covers above; their
orders are prime to $p$. Conversely a tame finite cover of
$\mathbf A^1$ would extend to a tame cover of $\mathbf P^1$
branched at only one point, contradicting Riemann--Hurwitz.
Finally, the equivalence between finite extensions of the excellent
henselian local fraction field and of its completion follows by
approximating the coefficients of a separable defining polynomial
and Hensel's lemma. Thus the same conclusion holds for the local
groups used in B.1, not only for completed fields.
\end{proof}

\paragraph{B.7 (Proved): complete comparison and its exact exceptions.}\label{B7}
\textit{Manuscript: LM \S14, folio 63, scans 129.}
Let $X/k$ be a smooth geometrically connected separated curve of
finite type, $S\subset X$ finite and reduced, and $F=(M,N_s,a_s)$
an arbitrary torsion abelian object of the full \`etale mixed
topos. The comparison $R\Gamma(\mathcal B,F)\to R\Gamma(X,f^*F)$
is an equivalence for every $F$, except for the following geometric
types:
\begin{enumerate}[label=(\roman*)]
\item $X$ is proper of genus zero and $S=\varnothing$;
\item $\operatorname{char}k=0$, $X\simeq\mathbf P^1_k$, and
$S$ consists of one geometric point (necessarily a $k$-point).
\end{enumerate}
In (i), set $M^{(p')}$ equal to the prime-to-$p$ part of $M$,
and to $M$ in characteristic zero. The comparison cone is
\begin{equation}\label{eq:B-genuszero}
 R\Gamma(k,M^{(p')}(-1))[-2].
\end{equation}
In (ii) the cone is $R\Gamma(k,M(-1))[-2]$, independently of
the boundary module $N_s$ and its specialization map. Thus these
formulas decide the comparison for \emph{each} coefficient object,
including arithmetic coefficients with acyclic Galois cohomology.
Over an algebraically closed field the first failure for every
nonzero indicated $M$ is exactly degree two. For $p$-primary
coefficients in characteristic $p$ there is no exception.

\begin{proof}
When $S$ is nonempty, $U$ is affine. B.5 gives its comparison and
B.6 gives injectivity of the geometric inertia maps except in
(ii). Over an arbitrary field, purely inseparable invariance first
allows passage to the perfect closure, taking the reduced boundary
after base change. This leaves the etale topoi and local finite-cover
categories unchanged. The residue fields are then finite separable
extensions of the ground field. The kernel of $(r_s,q_s)$ is
geometric; the residue map is injective into the ground-field
Galois group after fixing a residue-field embedding. Thus the
geometric result makes $(r_s,q_s)$ injective, and B.3 applies.
Alternatively, the geometric comparisons descend by the
compatible Hochschild--Serre spectral sequences. When $S$ is
empty, B.5 applies except in (i).

In (i), $\pi_1(X)=G_k$ and $f$ is the structural morphism.
On a geometric genus-zero curve, $R^0f_*f^*M=M$,
$R^1f_*f^*M=0$, $R^2f_*f^*M=M^{(p')}(-1)$, and the remaining
terms vanish. These are the Kummer and Artin--Schreier calculations
in B.5. The cone of the unit $M\to Rf_*f^*M$ is therefore
$M^{(p')}(-1)[-2]$; apply $R\Gamma(k,-)$. This conclusion does
not require the genus-zero curve to split, nor does it assert a
splitting of its derived direct image.

In (ii), $\mathcal B$ is the arrow topos over $\beta G_k$.
Geometric calculation B.4 gives its relative comparison cone
$M(-1)[-2]$ with the natural Galois action. Descent gives the
stated cone. In every other case the cone is zero. The fieldwise
formulation avoids a false uniform assertion about the first
arithmetic degree: if $X$ is a nonsplit real conic and $M=\mathbb
Z/2$, its arithmetic cohomology is that of $\mathbf{RP}^2$,
whereas $H^*(G_{\mathbb R},\mathbb Z/2)=\mathbb F_2[t]$.
The comparison first fails in degree three. For a split genus-zero
curve a rational point splits the unit and the degree-two
orientation class already witnesses failure. More generally
\eqref{eq:B-genuszero} and its long exact sequence specify the
failure in every arithmetic degree without a hidden splitting
assumption.
For the real example, the complex points of the anisotropic
conic identify with $S^2$, with complex conjugation the free
antipodal involution (in a projective coordinate it is
$z\mapsto-1/\overline z$). Etale--Betti comparison over
$\mathbb C$ and descent along $X_{\mathbb C}\to X$ identify
the arithmetic complex with the equivariant cohomology
of that sphere. Since the action is free, this is the
cohomology of $\mathbf{RP}^2$. The unit is the classifying
map of its double cover; it is an isomorphism in degrees
zero, one, and two, while degree three of the target
curve vanishes and degree three of $BC_2$ is nonzero.
\end{proof}

\paragraph{B.8 (Refuted, with replacement proved): tame groups for $p$-torsion.}\label{B8}
\textit{Manuscript: LM \S12 and 14, folio 53 and 63, scans 109 and 129.}
The prime-to-$p$ version of B.7 is obtained using tame inertia
and tame local systems, with the additional exceptional geometric
pair $(X_{\overline k},(S_{\overline k})_{\mathrm{red}})
\simeq(\mathbf P^1_{\overline k},\{\infty\})$ in
\emph{every} characteristic. Over an imperfect field this includes
a nonsplit genus-zero curve with one purely inseparable boundary
point; the perfection reduction in B.7 covers it.
It cannot be extended to $p$-torsion by retaining only tame groups.
For $U=\mathbf A^1_k$, with $k$ algebraically closed of
characteristic $p$, its tame group is
trivial but
\[
 H^1_{\et}(U,\mathbb F_p)=k[x]/\{h^p-h:h\in k[x]\}\ne0.
\]
The full-group statement B.7 remains valid for $p$-torsion.

\begin{proof}
For a tame finite local system of order prime to $p$, a
degree-one torsor trivializing a class is again tame: wild
inertia acts trivially on the coefficients and has no nontrivial
homomorphism to their prime-to-$p$ additive group. Together with
the curve cohomological-dimension argument, this proves the
tame $K(\pi,1)$ comparison. B.6 supplies tame local injectivity
outside the one-cusp rational case, and B.3 completes the proof.
For the counterexample, Artin--Schreier gives the displayed
cokernel. The class of $x$ is nonzero: a nonconstant polynomial
$h$ has $\deg(h^p-h)=p\deg h$, while a constant gives a constant.
This also proves that one cannot erase wild monodromy on the
grounds that the underlying variety is a curve. In particular,
the demand for an example of failure of the \emph{full} mixed
comparison with $p$-torsion on these smooth curves has a false
premise; the failure belongs to its tame replacement.
\end{proof}

\subsection{Comparison with logarithmic topoi}

\paragraph{B.9 (Proved): the logarithmic comparison, including coefficients.}\label{B9}
\textit{Manuscript: LM \S14, folio 60--62, scans 123--127.}
Give the smooth curve $X$ the divisorial log structure of $S$.
There are canonical geometric morphisms
\[
 X_{\mathrm{ket}}\xrightarrow{\epsilon}X_{\et}
       \xrightarrow{f}\mathcal B_{X;S}.
\]
This does not identify any two of these topoi. On the closed
fibre over $s$, the logarithmic topos is $\beta H_s^t$, where
$H_s^t=H_s/P_s$ and $P_s$ is wild inertia. The mixed topos has
closed fibre $\beta Q_s$. Its comparison sends a $Q_s$-set to
its inflation to $H_s^t$. In characteristic zero $P_s=1$.

More explicitly, the glued covering model of logarithmic sheaves
which are essentially locally constant on $U$ consists of
\[
 (M,V_s,b_s),\qquad
 V_s\in\beta H_s^t,\qquad
 b_s:(H_s\to H_s^t)^*V_s\longrightarrow r_s^*M.
\]
Here ``covering model'' means the associated glued covering
topos with its comparison to the full logarithmic topos, not
that every full logarithmic sheaf is such an object.
The inverse image of $(M,N_s,a_s)$ is
$(M,(H_s^t\to Q_s)^*N_s,a_s)$. Thus its boundary monodromy
is trivial on tame inertia. A locally constant logarithmic
sheaf can have nontrivial boundary inertia, so is in general
outside this image.

For $\Lambda=\mathbb Z/n$ with $n$ invertible on $X$, and every
torsion $\Lambda$-module object $F$ of $\mathcal B_{X;S}$, put
$\mathcal F=f^*F$ and $N_s=i_s^*\mathcal F$. There is the exact
triangle
\[
 R\Gamma(X,\mathcal F)\longrightarrow
 R\Gamma(X_{\mathrm{ket}},\epsilon^*\mathcal F)
 \longrightarrow
 \bigoplus_{s\in S}R\Gamma(k(s),N_s(-1))[-1]
 \longrightarrow R\Gamma(X,\mathcal F)[1].                 \tag{B.9.1}
\]
The boundary is the logarithmic residue/Gysin boundary; the
triangle need not split. For a finite locally constant
logarithmic $\Lambda$-sheaf $L$ one instead has
\[
 R\Gamma(X_{\mathrm{ket}},L)\simeq R\Gamma(U,L|_U).
                                                               \tag{B.9.2}
\]
These assertions do not extend to $p$-torsion by keeping the
same Tate-twist formula.

\begin{proof}
The morphism $\epsilon$ is induced by viewing an ordinary
etale morphism as a strict Kummer-etale morphism. A strict
local trait acquires, through all Kummer roots of a
uniformizer, the tame inertia group
$\prod_{\ell\ne p}\mathbb Z_\ell(1)$. Its closed logarithmic
point therefore has group $H_s^t$. A section of a nearby
$H_s$-set which survives this system is precisely a
$P_s$-invariant section. Hence its gluing functor is
$(r_s^*M)^{P_s}$, now with its $H_s^t$-action. Adjunction gives
the displayed triples. This also proves the inflation formula
for the comparison, without conflating wild and tame
monodromy. A nontrivial Kummer character of tame inertia gives
the promised object outside its image.

Every point of the mixed topos lies in its open part or one
of its closed parts: pull its open subterminal back to
$\Set$, where it is either $0$ or $1$. Over an algebraically
closed field its closed point has no automorphism, whereas
the fibre functor on the logarithmic point has automorphism
group $\widehat{\mathbb Z}^{(p')}(1)$. In the mixed topos all
points on the connected open stratum are identified at the
level of its covering theory. Full $X_{\et}$ and
$X_{\mathrm{ket}}$ still distinguish different underlying
locations by their open subterminals. For instance the mixed
topos has only $1+2^{|S|}$ subterminals in this connected
geometric situation, whereas the full topoi retain every
Zariski open of $X$. These are object-level and point-level
distinctions independent of cohomology.

For the cohomological assertion, on a strict local trait the
continuous cohomology of tame inertia with constant
prime-to-$p$ coefficients is $\Lambda$ in degree zero,
$\Lambda(-1)$ in degree one, and zero above. The Kummer
sequence makes this identification canonical. Off $S$ the
log structure is trivial. Consequently
\[
 \Lambda\longrightarrow R\epsilon_*\Lambda\longrightarrow
 i_*\Lambda(-1)[-1]\longrightarrow\Lambda[1],
\]
with $i:S\hookrightarrow X$. The projection formula gives
\eqref{B9-proj} below; it can also be checked on these strict
local stalks, whose two cohomology modules are free over
$\Lambda$:
\begin{equation}\label{B9-proj}
 R\epsilon_*\epsilon^*\mathcal F
     \simeq\mathcal F\otimes_\Lambda^{\mathbf L}R\epsilon_*\Lambda.
\end{equation}
The local calculation and projection formula are the
curve case of \citet[Theorem~5.2 and \S5.4]{Illusie2002}.
Tensoring the triangle and taking sections proves (B.9.1).
For constructible coefficients this is the usual derived
projection formula; arbitrary torsion coefficients follow
by filtered unions of finite subobjects and continuity.

For (B.9.2), a logarithmically locally constant sheaf is
trivialized by a Kummer-etale covering. On such a covering
the assertion $L\simeq Rj_{\mathrm{ket}*}L|_U$ reduces to
$\Lambda\simeq Rj_{\mathrm{ket}*}\Lambda$. The latter is
logarithmic purity, \citet[Theorem~7.4]{Illusie2002}; in this
curve case its proof is the same Kummer calculation of
$R^1\epsilon_*\Lambda$ and $R^1j_*\Lambda$, with higher
stalks zero. Descent and global sections give (B.9.2).
Over $\mathbb C$, Theorem~5.9 of that source identifies the
same computation with the real oriented blow-up: its
boundary circles retain the local loops. Finally B.8's
Artin--Schreier class proves the asserted characteristic-$p$
limitation. For example on $(\mathbb P^1,\infty)$ in
characteristic zero, constant coefficients have mixed
$H^2=0$, ordinary etale $H^2=\Lambda(-1)$, and logarithmic
$H^2=0$. Agreement of two numbers in this example does not
give an equivalence of topoi.
\end{proof}

\subsection{Orientation, support, and trace}

\paragraph{B.10 (Refuted, with replacement): the exact orientation formulas.}\label{B10}
\textit{Manuscript: LM \S14, folio 63--64, scans 130--131.}
In this paragraph the groups are geometric, over an
algebraically closed field. In characteristic zero put
$\widehat{\mathbb Z}^{\dagger}=\widehat{\mathbb Z}$ and
$T=\widehat{\mathbb Z}(1)$; in characteristic $p>0$ omit the
$p$-factor and write $T=\widehat{\mathbb Z}^{(p')}(1)$.
Let $\overline G=\pi_1(X)$ for a smooth proper connected
curve $X$, and $G=\pi_1(U)$. For $g(X)\geq1$ there are
canonical isomorphisms
\[
 H^2(\overline G,T)\simeq\widehat{\mathbb Z}^{\dagger},
 \qquad H^2(\overline G,\widehat{\mathbb Z}^{\dagger})
                         \simeq T^\vee.                    \tag{B.10.1}
\]
If a boundary inertia is cohomologically effective as in
B.3, its support group is
\[
 H^2_s(\mathcal B_{X;S},\widehat{\mathbb Z}^{\dagger})
       \simeq L_s^\vee,\qquad
       L_s=\widehat{\mathbb Z}^{\dagger}(1).                \tag{B.10.2}
\]
For $g\geq1$ the local-to-global map identifies all the
$L_s$ canonically with the dual of the second group in
(B.10.1). If $S\ne\varnothing$, the unfilled group instead
has
$H^2(G,\widehat{\mathbb Z}^{\dagger})=0$.
Thus scan~131 must use the \emph{filled} group for its
genus-positive $H^2$ formula. For genus zero and no puncture,
the displayed group-cohomology formula is false.

For every nonempty $S$, including a single puncture on
$\mathbb P^1$, define the relative support complex using the
\emph{full} peripheral homomorphisms by
\[
 R\Gamma_c(G,\{H_s\},\Lambda)
 :=\operatorname{fib}\left(R\Gamma(G,\Lambda)
                  \longrightarrow\bigoplus_sR\Gamma(H_s,\Lambda)\right).
                                                               \tag{B.10.3}
\]
Then
\[
 H^2_c(G,\{H_s\},\widehat{\mathbb Z}^{\dagger})\simeq T^\vee.
                                                               \tag{B.10.4}
\]
This equals $H^2(\mathcal B,j_!\widehat{\mathbb Z}^{\dagger})$
when the local maps are effective. For
$(\mathbb P^1,\infty)$ in characteristic zero the latter
group is zero, so the unqualified marginal claim for the
ordinary mixed topos is refuted.

\begin{proof}
First work with $\Lambda=\mathbb Z/n$, $n$ prime to the
characteristic. The Kummer sequence and the degree map give
$H^2(X,\Lambda(1))\simeq\Lambda$; for the precise curve
calculation see \citet[Lemma~59.69.1]{StacksCurveKummer}.
The $K(\pi,1)$ result B.5 gives (B.10.1) in positive genus.
For a strict local trait, the cone of
$\Lambda\to R\Gamma(I_s,\Lambda)$ is
$\operatorname{Hom}(L_s,\Lambda)[-1]$. The support fibre is
its shift by $-1$, proving (B.10.2) when B.3 identifies the
effective and full local groups. The boundary class of a
uniformizer is the Kummer residue, and its Gysin image is
the degree-one point class. Its trace is $1$. Therefore the
identifications for different $s$ agree after the
local-to-global support map. This proves the claimed
independence without choosing generators of inertia.

Affineness and B.5 imply the vanishing of $H^2(G,\Lambda)$.
Localization on the proper curve gives
$R\Gamma(X,j_!\Lambda)$ as the fibre in (B.10.3), because
the geometric local punctures are exactly the $H_s$ and
$U$ is a $K(\pi,1)$. Its top compact-support trace is
$\Lambda(-1)$, proving (B.10.4). In the exceptional affine
line case the fibre is explicitly
\[
 \operatorname{fib}\bigl(\Lambda\to
 R\Gamma(\widehat{\mathbb Z}(1),\Lambda)\bigr)
                  \simeq\Lambda(-1)[-2],
\]
whereas B.4 proves the mixed $j_!$ group is zero. A proper
geometric genus-zero curve has trivial fundamental group,
so its positive group cohomology vanishes. Finally take the
derived inverse limit in $n$. The finite groups in these
formulas have the surjective transition maps of the stated
Tate module, so no unwanted $\varprojlim^1$ term occurs in
these degrees. This explains the meaning of the completed
coefficients. In characteristic $p$, $H^2(X,\mathbb F_p)=0$
by Artin--Schreier, hence no all-primes rank-one orientation
formula can replace the stated prime-to-$p$ formula.
\end{proof}

\paragraph{B.11 (Proved): the cohomological trace.}\label{B11}
\textit{Manuscript: LM \S14, folio 64, scans 131--132.}
Let $G'\to G$ be a continuous homomorphism of geometric
curve groups with the peripheral data of scans~131--132:
the loops in $S'_0$ map to specified target loops with the
specified degrees, and the loops in $S'_1$ map trivially.
Include the stated identification $T'\simeq T$ of
orientation modules, denoted $\phi:T'\xrightarrow{\sim}T$,
and let $\Lambda=\mathbb Z/n$ with $n$ invertible in the field.
In every case where the group and
peripheral complexes compute the curve, there is a
functorial trace
\[
 \operatorname{Tr}_f:R\Gamma_c(U',\Lambda)
                      \longrightarrow R\Gamma_c(U,\Lambda)
                                                               \tag{B.11.1}
\]
which on $H^2_c$ is $(\phi^{-1})^\vee:T'^\vee\to T^\vee$
modulo $n$, the contragredient orientation isomorphism.
The isomorphism $\phi$ is explicitly part of the datum
on LM \S14, folio~64, scan~131; scan~132 specifies its
contragredient. Here the cohomology is defined by (B.10.3)
for punctured curves, and by group cohomology for proper
curves of positive genus. If the ordinary mixed topoi
compute these complexes, this is also their support
cohomology trace. Without that qualification the requested
mixed-topos trace assertion is false, by B.10.

\begin{proof}
We give a construction that works for an abstract
homomorphism and does not assume it comes from an
algebraic map. For a smooth geometric curve the cup product
and the trace of its orientation class give Poincare
duality
\[
 R\Gamma_c(U,\Lambda(1))\simeq
 R\operatorname{Hom}_\Lambda(R\Gamma(U,\Lambda),\Lambda)[-2].
                                                               \tag{B.11.2}
\]
For the etale statement and its reduction by specialization,
see \citet[VI, Theorem~2.2 and its proof]{DeligneArcata1977}.
In the present setting it can be proved by a finite cell
complex for an oriented surface with boundary, followed
by the prime-to-characteristic comparison, or by the
Kummer local calculation and curve duality. The finite
complex has one dual cell for each cell, and intersection
of a cell with its dual is $1$; its boundary matrices are
transposes with the usual signs. This proves the perfect
pairing, including the boundary case. The etale version
has the same trace normalization. All the constant
cohomology modules here are finite free over $\Lambda$,
so these are perfect complexes.

Using B.5, the homomorphism induces restriction
$R\Gamma(G,\Lambda)\to R\Gamma(G',\Lambda)$.
Apply $R\operatorname{Hom}_\Lambda(-,\Lambda)[-2]$ and the
two isomorphisms (B.11.2), and then remove the twist by
the given $T'\simeq T$. This is (B.11.1). Composition is
functorial because restriction and dualization are. In
degree two the map is dual to
$H^0(U,\Lambda)\to H^0(U',\Lambda)$, which is the identity
on $\Lambda$ for connected curves. Hence it is precisely
the orientation isomorphism, rather than multiplication
by a global covering degree. For a finite covering the
usual relation $\operatorname{Tr}_f f^*=\deg(f)$ still
holds: pullback, not this top trace, multiplies the
fundamental cohomology class by the degree. The construction
is equivariant under operators preserving the supplied
orientation data, and therefore descends through
Hochschild--Serre in arithmetic situations. No general
six-functor formalism for arbitrary topoi is asserted.
\end{proof}

\subsection{Mixed maps and operators}

\paragraph{B.12 (Proved): the mixed morphism, with its commutation data.}\label{B12}
\textit{Manuscript: LM \S14, folio 64--67, scans 132--137.}
Let $S'=S'_0\amalg S'_1$. To construct
$g:\mathcal B'\to\mathcal B$, give the following coherent
data: a geometric morphism $\beta G'\to\beta G$;
for $t\in S'_0$ a target $\tau(t)$ and maps of local and
residue groupoids with invertible commutation two-cells;
and, for $t\in S'_1$, a geometric morphism
$p_t:\beta Q'_t\to\beta G$ together with an invertible
two-cell
\[
 (\beta H'_t\to\beta G'\to\beta G)
       \simeq(\beta H'_t\to\beta Q'_t\xrightarrow{p_t}\beta G).
                                                               \tag{B.12.1}
\]
The inverse image on triples is
\[
 M'=f^*M,\qquad
 N'_t=\begin{cases}
        h_t^*N_{\tau(t)},&t\in S'_0,\\
        p_t^*M,&t\in S'_1.
       \end{cases}                                           \tag{B.12.2}
\]
For the first case use the transported $a_{\tau(t)}$;
for the second use (B.12.1) itself. These formulas
construct the desired geometric morphism and its open
and filled comparisons. In geometric discrete notation
$p_t$ is the right $G$-torsor $P_t$ of scan~134.
For an infinite profinite $G$ the word torsor here means
the compatible system of finite torsors, or the
equivalent profinite/localic torsor, not a discrete
regular $G$-set.

\begin{proof}
An object of the target is the diagram in B.1. Applying
the specified inverse images to this diagram and using
the commutation isomorphisms gives exactly (B.12.2), with
the required equivariance. The same prescription on
componentwise maps is a functor. In the original triple
description finite limits and all colimits are computed
on $M$ and the $N_s$: the functors $q_s^*$ and $r_s^*$
preserve both. Each functor in (B.12.2) is an inverse
image and has the same preservation properties.
Consequently the constructed functor is left exact and
cocontinuous, and, by accessibility, has a right adjoint.
It is the inverse image of a geometric morphism.
Composition of the formulas gives the composition of
the maps, with exactly the composition of their
two-cells. Restriction to the open part gives $f$;
the filled comparison follows by quotienting the
peripheral normal subgroups as in B.2.

For $Q'_t=1$ the functor $p_t^*$ is the fibre functor
$M\mapsto P_t\times^G M$. Equation (B.12.1) says that
the local inertia acts trivially on that fibre functor,
with a specified identification with the restriction
from the open stratum. Merely listing an isomorphism
class of torsors and omitting this datum would not
specify a mixed map. This is exactly the isocommutative
triangle and its collapsed-loop counterpart on scan~134.
For profinite groups, replacing $P_t$ by its finite
quotients proves the same formula on every continuous
set. A free discrete $G$-torsor would be discontinuous
when $G$ is infinite, which explains the qualification.
\end{proof}

\paragraph{B.13 (Proved): operators, splittings, and the exact rigidity claim.}\label{B13}
\textit{Manuscript: LM \S14, folio 67, scans 136--137.}
Suppose a group $\Gamma$ acts coherently on the preceding
data. On an orbit of $S'_1$ choose $t$ with stabilizer
$\Gamma_t$. Let
\[
 1\longrightarrow G\longrightarrow E\longrightarrow\Gamma
       \longrightarrow1
\]
represent the action on the connected covering groupoid.
The datum of an equivariant point over that orbit is a
right $G$-torsor $P$ together with a splitting
$q:\Gamma_t\to E^P$ of the twisted extension, over
$\Gamma_t$. Its automorphism group is
\[
 C_{G^P}(q(\Gamma_t))=(G^P)^{\Gamma_t}.                    \tag{B.13.1}
\]
If this group is trivial, every two isomorphic such
objects have exactly one isomorphism. This proves the
rigidity sentence of scan~136. It neither asserts that
a splitting exists nor that different splitting classes
are isomorphic.

\begin{proof}
An equivariant structure on the fibre torsor is a
compatible lift of every element of $\Gamma_t$ to an
automorphism over its prescribed action. Such lifts
form the extension $E^P$; the compatibility for
products is exactly the homomorphism condition for
a section. Equivalently, after choosing a point of
$P$, it is a lift into $E$, and changing the point
conjugates the section by $G$. This proves both the
torsor-and-section description and its independence
of the choice. An automorphism $a$ of $P$ preserves
the equivariant structure precisely when
$a q(\gamma)a^{-1}=q(\gamma)$ for every $\gamma$.
This is (B.13.1). The set of isomorphisms between two
isomorphic objects is a torsor under the automorphism
group of either, proving unique isomorphism when that
group is trivial. In the profinite version require
all actions and splittings continuous and interpret
the torsor as in B.12; the identical equations hold
at every finite quotient and pass to the inverse
system. Extending over all orbits and using B.12
constructs the equivariant mixed morphism. The
subsequent rational-point conjecture on scan~137 is
not a consequence of this elementary rigidity
calculation.
\end{proof}
